\documentclass[11pt]{article}
\usepackage[utf8]{inputenc}
\usepackage{amsmath,amssymb}
\usepackage{authblk}
\usepackage[margin=1in]{geometry}
\usepackage[hidelinks]{hyperref}
\usepackage{xurl}
\usepackage{setspace}

\title{Charge Densities in Real Space:\\
Primitive Ontology and the Individuation of Electrons}
\author[1]{Claes Johnson}
\affil[1]{KTH Royal Institute of Technology, SE-100 44 Stockholm, Sweden \\
\texttt{claesjohnson@gmail.com}}
\date{\today}

\begin{document}
\maketitle

\begin{abstract}
We compare the ontology of an atomic system of $N$ electrons under two models. Standard quantum
mechanics treats the electrons as \emph{identical particles} and represents the system by a probability
amplitude over their configurations in a $3N$-dimensional configuration space, antisymmetric under
exchange. \emph{RealQM} represents
it by $N$ unit charge densities, each carried on its own region $\Omega_i$ of ordinary three-dimensional
space, the regions non-overlapping ($\Omega_i\cap\Omega_j=\varnothing$). Two long-standing
metaphysical problems --- the \emph{arena} of the wave function and the \emph{individuality} of identical
particles --- attach to the first model, one to each of its two features: the $3N$-dimensional domain and
the antisymmetry. The second model has neither feature and so incurs neither problem. It is a primitive
ontology in three-space (no configuration-space object exists), and it \emph{individuates} its electrons
absolutely, by the regions they occupy (the Principle of the Identity of Indiscernibles is satisfied); the
two virtues are one fact, the non-overlap of the $\Omega_i$ doing double duty. RealQM is
moreover \emph{constructive}: its charge densities are the objects one actually computes, applied to real
chemistry and nuclear structure. We state both ontologies in equations, place RealQM against the Standard
Model's point-and-field ontology for perspective, and bound the claim: RealQM is a rival theory,
empirically warranted in its atomic and nuclear domain, and the metaphysical verdict is conditional on
that adequacy.
\end{abstract}

\setstretch{1.05}

\section{Two problems in one formula}
\label{sec:two}

An atomic system of $N$ electrons admits two descriptions. Standard quantum mechanics treats the electrons
as \emph{identical particles} and represents the system by one probability amplitude over their joint
configurations, in a space of $3N$ dimensions. RealQM treats them as $N$ \emph{individual charge
densities}, each carried on its own region $\Omega_i\subset\mathbb R^3$, in ordinary three-dimensional
space (\S\ref{sec:po}). This contrast --- identical particles in a high-dimensional configuration space
against individuated densities in ordinary space --- is the whole of what follows. We take the standard
side first.

Write its state:
\begin{equation}
\Psi:\ \mathbb R^{3N}\to\mathbb C,\qquad
\Psi(\dots,\mathbf x_i,\dots,\mathbf x_j,\dots)=-\,\Psi(\dots,\mathbf x_j,\dots,\mathbf x_i,\dots).
\label{eq:std}
\end{equation}
The domain is $\mathbb R^{3N}$; the minus sign is the \emph{antisymmetry} of the amplitude under exchange of
any two particles --- Fermi--Dirac statistics, the Pauli principle in its formal guise. Two metaphysical
difficulties follow directly, one from each of these two features of \eqref{eq:std}.

\emph{The arena}, from the domain $\mathbb R^{3N}$. If $\Psi$ is the fundamental object --- wave-function
realism \cite{Albert1996,NeyAlbert2013,Ney2021} --- then ordinary three-space is not fundamental but must
be recovered from a field on a $3N$-dimensional space. The rival \emph{primitive ontology} programme
\cite{Maudlin2007,Allori2013,DGZ1992} insists instead on beables in $\mathbb R^3$, with $\Psi$ demoted to
a dynamical or nomological role. Either way the $3N$-dimensional domain is the source of the difficulty.

\emph{Individuality}, from the antisymmetry. The permutation symmetry in \eqref{eq:std} makes the
particles, on the received reading, \emph{non-individuals}: no fact settles which is which, and two
electrons in one atom are numerically distinct yet share all properties, in apparent violation of
Leibniz's Principle of the Identity of Indiscernibles (PII) \cite{FrenchKrause2006,FrenchSEP}. Defenders of
PII retreat to \emph{weak discernibility} by irreflexive relations, with no distinguishing monadic
property \cite{Saunders2006,MullerSaunders2008}.

Both problems are consequences of the two choices in \eqref{eq:std}: the domain $\mathbb R^{3N}$, and the
antisymmetry. A formulation making neither choice does not incur them.

\section{A real-space ontology, in equations}
\label{sec:po}

Take instead the RealQM formulation \cite{RealQM}
\begin{equation}
\Psi(\mathbf x)=\sum_{i=1}^N\psi_i(\mathbf x),\quad \mathbf x\in\mathbb R^3,\qquad
\operatorname{supp}\psi_i=\overline{\Omega_i},\quad \Omega_i\cap\Omega_j=\varnothing\ (i\neq j),
\label{eq:real}
\end{equation}
where $\psi_i^2$ is the unit charge density of electron $i$ on its own region $\Omega_i\subset\mathbb R^3$
and $\rho=\sum_i\psi_i^2$ is the total electron density. For a nucleus (kernel) of charge $Z$ at the
origin, with attractive potential $V(\mathbf x)=Z/|\mathbf x|$, the ground state minimises the energy
\begin{equation}
E=\sum_i\tfrac12\!\int_{\Omega_i}|\nabla\psi_i|^2
-\int V\rho\,dx
+\tfrac12\!\iint\frac{\rho(x)\rho(y)}{|x-y|}\,dx\,dy,\qquad
\int_{\Omega_i}\psi_i^2\,dx=1,
\label{eq:E}
\end{equation}
over the $\psi_i$ and the free boundaries $\partial\Omega_i$.

RealQM is, at once, a \emph{primitive ontology in three-space}: what exists is the extended charge density
$\rho$ in $\mathbb R^3$, and there is \emph{no} object on $\mathbb R^{3N}$ to interpret or to recover
three-space from --- \eqref{eq:std}'s domain is simply absent. Unlike primitive-ontology theories that
carry a matter density \emph{alongside} a configuration-space $\Psi$ \cite{Allori2013}, here the density is
the whole matter ontology and \eqref{eq:E} its law. It is a rival theory, not a reinterpretation of
\eqref{eq:std}, warranted by the chemistry and nuclear structure it computes \cite{RealQM}.

\section{Individuation by occupation}
\label{sec:ind}

The basic objects of \eqref{eq:real} are extended and disjoint, and this individuates them. Define electron
$i$ as the density on $\Omega_i$; then occupancy of $\Omega_i$ is a \emph{monadic} property that $i$ has
and $j\neq i$ lacks, since $\Omega_i\cap\Omega_j=\varnothing$. Hence
\begin{equation}
i\neq j\ \Longleftrightarrow\ \Omega_i\neq\Omega_j:\qquad
\text{no two electrons share all properties.}
\end{equation}
The Principle of the Identity of Indiscernibles, standardly taken to fail for \eqref{eq:std}, is
\emph{satisfied}. Moreover the electrons are \emph{absolutely} discernible --- each has a distinguishing
monadic property (its region) --- which is stronger than the \emph{weak} discernibility
\cite{MullerSaunders2008} to which PII's defenders retreat under \eqref{eq:std}. Individuality in the full
classical sense is recovered.

And extension and individuation are \emph{the same fact}. The disjointness $\Omega_i\cap\Omega_j=\varnothing$
that keeps each $\psi_i$ a density in $\mathbb R^3$ (the primitive-ontology virtue) is the very disjointness
that discerns the electrons (the individuality virtue). Where \eqref{eq:std} encodes exclusion as
antisymmetry --- a fact about a configuration-space function --- \eqref{eq:real} encodes it as
impenetrability --- a fact about \emph{space}: two charge clouds cannot occupy one place. Pauli exclusion
and the identity of indiscernibles, distinct in the received picture, are one geometric fact here.

\section{Computability: the ontology is the method}
\label{sec:comp}

The standard amplitude \eqref{eq:std} lives on $\mathbb R^{3N}$, and representing it there costs
exponentially in the number of particles: solving \eqref{eq:std} directly is feasible only for very small
$N$, and beyond a few electrons one must retreat to approximations (Hartree--Fock, density-functional
theory). RealQM \eqref{eq:real}--\eqref{eq:E} lives instead on ordinary $\mathbb R^3$, where complexity
scales with mesh points in three dimensions, not with the exponential dimension of $\mathbb R^{3N}$; it is
directly computable for real systems.

This also sets RealQM apart from the metaphysical proposals of \S\ref{sec:two}. Wave-function realism and
the primitive-ontology programme are readings of a formalism one solves --- approximately --- on
$\mathbb R^{3N}$; the ontology is interpretive, supplied after the calculation. In
\eqref{eq:real}--\eqref{eq:E} the beables \emph{are} the objects of computation: one minimises the energy
\eqref{eq:E} directly over the charge densities $\psi_i$ in $\mathbb R^3$, and the primitive ontology
carries the calculation rather than glossing it. The
theory is moreover \emph{applied}: the same functional \eqref{eq:E} returns atomic energies across the
periodic table, molecular geometries and bond energies, reaction paths, condensed phases, and --- in a
proton--electron reading --- nuclear binding. These are not paper derivations but a large public library of
interactive, browser-runnable simulations, collected in the RealQM gallery \cite{RealQMgallery}, where the
charge densities of atoms, molecules, reactions, and nuclei can be computed and watched directly. The
metaphysical significance is that RealQM's ontology is not posited for interpretive
tidiness but \emph{earns its keep by doing the calculation}: it is a working, applied method of which the
ontology of \S\S\ref{sec:po}--\ref{sec:ind} is the direct reading. An ontology that is also the algorithm
is a stronger kind of proposal than one available only as an interpretation.

\section{The point and the field: a perspective}
\label{sec:sm}

The extension at issue reaches past non-relativistic quantum mechanics. The Standard Model's Lagrangian
carries three gauge forces and a Higgs and Yukawa sector,
\begin{equation}
\mathcal L_{\rm SM}=-\tfrac14 F^a_{\mu\nu}F^{a\mu\nu}
+\bar\psi\,i\gamma^\mu D_\mu\psi+|D_\mu\Phi|^2-V(\Phi)-\big(y\,\bar\psi_L\Phi\psi_R+\text{h.c.}\big),
\end{equation}
its fundamental carriers \emph{point} excitations of quantum fields. That is a second departure from
\emph{res extensa}: where \eqref{eq:std} moved the object into $\mathbb R^{3N}$, field theory shrank it to
a point in $\mathbb R^3$. The point carries the classical infinity of a zero-extension charge's
self-energy, softened but not cured by renormalisation --- and it is telling that string theory restores
extension (to strings) precisely to tame it. Retaining only the electromagnetic term of $\mathcal L_{\rm
SM}$ and reading it, non-relativistically and electrostatically, as extended densities returns exactly the
energy functional \eqref{eq:E}: RealQM is that one term, read as \emph{res extensa} in $\mathbb R^3$. The
ontological choice --- the point or the extended thing --- is thus one question asked at every scale, and
RealQM answers it, in its domain, on the side of extension. Nothing here
disputes the Standard Model's account of the sub-atomic world; the scales differ.

\section{Conclusion}
\label{sec:concl}

The arena of the wave function and the individuality of particles are two readings of the two features of
one formula, \eqref{eq:std}: its $3N$-dimensional domain and its antisymmetry. RealQM, the real-space
charge-density formulation \eqref{eq:real}--\eqref{eq:E}, declines both, and so meets both problems with one
move --- a primitive ontology in three-space whose electrons are individuated absolutely, by the regions
they occupy, because $\Omega_i\cap\Omega_j=\varnothing$ does double duty. Extension and individuation are
the same fact, and Pauli exclusion and the identity of indiscernibles coincide. Whether RealQM is
empirically adequate is for physics to decide; if it is, its metaphysics is the cleaner one.

\begin{thebibliography}{99}
\bibitem{Albert1996} D.~Z.~Albert, \emph{Elementary Quantum Metaphysics}, in J.~Cushing, A.~Fine, and
S.~Goldstein (eds.), \emph{Bohmian Mechanics and Quantum Theory: An Appraisal} (Kluwer, 1996), 277--284.
\bibitem{NeyAlbert2013} A.~Ney and D.~Z.~Albert (eds.), \emph{The Wave Function: Essays on the Metaphysics
of Quantum Mechanics} (Oxford University Press, 2013).
\bibitem{Ney2021} A.~Ney, \emph{The World in the Wave Function} (Oxford University Press, 2021).
\bibitem{Maudlin2007} T.~Maudlin, \emph{The Metaphysics Within Physics} (Oxford University Press, 2007).
\bibitem{Allori2013} V.~Allori, \emph{Primitive Ontology and the Structure of Fundamental Physical
Theories}, in \cite{NeyAlbert2013}, 58--75.
\bibitem{DGZ1992} D.~D\"urr, S.~Goldstein, and N.~Zangh\`i, \emph{Quantum Equilibrium and the Origin of
Absolute Uncertainty}, J. Stat. Phys. \textbf{67}, 843 (1992).
\bibitem{FrenchKrause2006} S.~French and D.~Krause, \emph{Identity in Physics: A Historical, Philosophical,
and Formal Analysis} (Oxford University Press, 2006).
\bibitem{FrenchSEP} S.~French, \emph{Identity and Individuality in Quantum Theory}, Stanford Encyclopedia
of Philosophy (2019).
\bibitem{Saunders2006} S.~Saunders, \emph{Are Quantum Particles Objects?}, Analysis \textbf{66}, 52 (2006).
\bibitem{MullerSaunders2008} F.~A.~Muller and S.~Saunders, \emph{Discerning Fermions}, British Journal for
the Philosophy of Science \textbf{59}, 499 (2008).
\bibitem{RealQM} C.~Johnson, \emph{Real Quantum Mechanics as Multiphase 3D Continuum Mechanics},
manuscript, 2026; \url{https://claes542.github.io/RealMolecule/gallery.html}.
\bibitem{RealQMgallery} C.~Johnson, \emph{The RealQM Gallery: interactive charge-density simulations of
atoms, molecules, reactions, and nuclei}, \url{https://claes542.github.io/RealMolecule/gallery.html}
(source: \url{https://github.com/Claes542/RealMolecule}).
\end{thebibliography}

\end{document}
